The inverse function theorem is often used in the following form, known as the implicit function theorem.

\begin{Theorem}{Implicit function theorem}
\label{thm:banach-implicit-function}
Let $X$, $Y$ and $Z$ be Banach spaces, and let $U$, $V$ be neighborhoods of the origin in $X$ and $Y$, respectively. Let $F$ be continuously differentiable from $U\times V$ into $Z$, and assume that $D_yF(0,0)\in\mathcal{L}(Y,Z)$ is one-to-one and onto. Assume, moreover, that $F(0,0)=0$. Then there exists a neighborhood $W$ of the origin in $X$ and a continuously differentiable mapping $f:W\to Y$ such that $F(x,f(x))=0$. Moreover, for small $x$ and $y$, $f(x)$ is the only solution $y$ of the equation $F(x,y)=0$.
\end{Theorem}

\begingroup
\renewcommand{\endProof}{\par\nobreak\hfill$\blacksquare$\par\vspace{0.1cm}}
\begin{Proof}
For the proof, we simply consider the mapping $\widetilde F:U\times V\to X\times Z$ defined by $\widetilde F(x,y)=(x,F(x,y))$ and apply the previous theorem to solve the equation $\widetilde F(x,y)=(x,0)$.
\end{Proof}
\endgroup

Although the implicit function theorem asserts the existence of a unique solution, it is often also useful in situations where uniqueness fails. Let us consider an equation $F(x,\lambda)=0$, where $F$ is a continuously differentiable function from $X\times\mathbb{R}$ to $Y$. We assume that $F(0,\lambda)=0$ for every $\lambda$. If $D_xF(0,\lambda)$ is one-to-one and onto, then the implicit function theorem tells us that $x=0$ is the only small solution of $F(x,\lambda)=0$. It is of interest to consider those values of $\lambda$ where this assumption fails. This leads to the subject of bifurcation theory, on which there is an extensive literature. We shall consider only the simplest case of a bifurcation. Specifically, we shall assume that $F$ is twice continuously differentiable, that $D_xF(0,0)$ has a one-dimensional nullspace spanned by $x_0$ and that its range has codimension one, and that $D_{x\lambda}F(0,0)x_0$ is not in the range of $D_xF(0,0)$. Let now $Z$ be a subspace of $X$ which complements the span of $x_0$; we substitute $x=\epsilon(x_0+z)$, where $z\in Z$. We then define $G(\epsilon,z,\lambda)=\frac{1}{\epsilon}F(\epsilon(x_0+z),\lambda)$, with the obvious definition $G(0,z,\lambda)=D_xF(0,\lambda)(x_0+z)$ in the limit $\epsilon=0$. Since $F$ was assumed of class $C^2$, we find that $G$ is still of class $C^1$. We can now use the implicit function theorem to solve the equation $G(\epsilon,z,\lambda)$ for $z$ and $\lambda$ as functions of $\epsilon$. To see this, we simply note that the derivatives are given by $D_zG(0,0,0)=\left.D_xF(0,0)\right|_Z$ and $D_\lambda G(0,0,0)=D_{x\lambda}(0,0)x_0$. The assumption that $D_{x\lambda}(0,0)x_0$ is not in the range of $D_xF(0,0)$ is precisely what is needed to guarantee the invertibility of the linearization. We thus obtain a bifurcating branch of nontrivial solutions of the equation $F(x,\lambda)=0$. This branch of nontrivial solutions is parametrized by the ``amplitude factor\textquotesingle\textquotesingle\ $\epsilon$.
